\subsection*{SOBREEXPOSICIÓN}
Investigar las causas\\
Estimar tiempos y distancias\\
Reportarlo a Seguridad Radiológica.

\subsection*{ACCIDENTES AUTOMOVILÍSTICOS}
Aplicar los tres criterios:
\begin{itemize}
    \item Garantizar la seguridad personal
    \item Asegurar integridad del contenedor
    \item Proteger al público
\end{itemize}

\subsection*{PÉRDIDA O ROBO DEL CONTENEDOR}
Avisar inmediatamente a Seguridad Radiológica.

\subsection*{CAÍDA DEL CONTENEDOR}
\begin{tabular}{ p{5.5cm} p{0.5cm} p{5.5cm} }
    Fuente afuera: & & Fuente dentro: \\
    \begin{minipage}[t]{5cm}
        \textbf{Si se puede introducir:} 
        \begin{itemize}
            \item Verificar carácter radiológico
            \item Llenar la forma SR-12
        \end{itemize}
    \end{minipage} 
    & 
    &
    \begin{minipage}[t]{5cm}
        \textbf{No se puede introducir:}
        \begin{itemize}
            \item Cortar la fuente
            \item Colocarla con blindaje especial
            \item Llenar forma SR-12
        \end{itemize}
    \end{minipage} \\
\end{tabular} 

\subsection*{FUENTE SUELTA}
\begin{itemize}
    \item \textbf{En manguera:} Sacarla al piso y llenar SR-12
    \item \textbf{En piso:} Colocarla en blindaje especial y llenar SR-12
\end{itemize}

Para una información más extensa se deberá consultar el Plan de Emergencias.

Para concluir, diremos que cualquier ahorro en protección radiológica aumenta considerablemente la probabilidad de que se produzca un accidente y de que las consecuencias sean lamentables.
